%% ============================================================================
%% Fichier  : annexes/Ann_A_glossaire.tex
%% Titre    : Glossaire
%% Version  : 1.0 -- novembre 2026
%% ============================================================================

\chapter{Glossaire}
\label{ann:glossaire}

Ce glossaire reprend, sous une forme abrégée, les notions définies dans les
chapitres. Le renvoi indique le chapitre où la notion est présentée.

\begin{description}
    \item[Action corrective] Action qui supprime la cause d'un écart pour éviter qu'il ne se reproduise~; la correction ne supprime que l'écart (Ch.~\ref{chap:rapport}).
    \item[Ampleur des pertes (LM)] Dans FAIR, pertes primaires et secondaires d'un événement (Ch.~\ref{chap:quantitatif}).
    \item[Analyse d'impact sur l'activité (BIA)] Analyse de l'impact d'une interruption au fil du temps, qui fixe priorités et objectifs de reprise (Ch.~\ref{chap:bia}).
    \item[Analyse de sensibilité] Étude de la variation d'un résultat quand on fait varier une hypothèse entre ses bornes (Ch.~\ref{chap:simulation}).
    \item[Appétence au risque] Types et niveau de risque qu'une organisation accepte pour atteindre sa mission (Ch.~\ref{chap:gouvernance}).
    \item[Appréciation du risque] Identification, analyse et évaluation du risque (ISO~31000) (Ch.~\ref{chap:cadres}).
    \item[Audit] Processus méthodique, indépendant et documenté d'obtention et d'évaluation de preuves au regard de critères (ISO~19011) (Ch.~\ref{chap:audit}).
    \item[Biais] Erreur systématique de jugement produite par une heuristique (Ch.~\ref{chap:perception}).
    \item[Cellule de crise] Équipe constituée à l'avance qui centralise l'information, décide et coordonne pendant la crise (Ch.~\ref{chap:cellule}).
    \item[Condition latente] Faiblesse créée par une décision d'organisation, qui laisse passer les défaillances actives (Reason) (Ch.~\ref{chap:retex}).
    \item[Constat d'audit] Résultat de la comparaison des preuves aux critères~: critère, fait, preuve, cause, conséquence, classement (Ch.~\ref{chap:conduire}).
    \item[Continuité d'activité] Capacité à fournir produits et services à un niveau acceptable pendant une perturbation (Ch.~\ref{chap:pca}).
    \item[Courbe de dépassement] Probabilité que la perte annuelle dépasse chaque montant (Ch.~\ref{chap:simulation}).
    \item[Crise] Situation instable, très incertaine, qui perturbe l'essentiel ou la crédibilité et exige une action urgente (ISO~22361) (Ch.~\ref{chap:crise}).
    \item[Critères d'audit] Référentiel auquel on compare les preuves~: norme, politique, contrat, loi (Ch.~\ref{chap:audit}).
    \item[Critères de risque] Termes de référence fixés à l'avance pour évaluer l'importance d'un risque (Ch.~\ref{chap:cadres}).
    \item[Culture juste] Culture où l'on signale erreurs et quasi-accidents sans crainte de sanction injuste (Ch.~\ref{chap:perception}).
    \item[Décision fondée sur la reconnaissance] Mode de décision des experts sous pression~: reconnaître, imaginer une action, la simuler (Klein) (Ch.~\ref{chap:decider}).
    \item[Déclaration d'applicabilité] Document d'un système de management de la sécurité qui justifie chaque mesure retenue ou exclue (Ch.~\ref{chap:ebios}).
    \item[Délai de reprise visé (RTO)] Durée au terme de laquelle l'activité doit avoir repris à un niveau minimal (Ch.~\ref{chap:bia}).
    \item[EBIOS Risk Manager] Méthode de l'ANSSI d'appréciation des risques par scénarios, en cinq ateliers (Ch.~\ref{chap:ebios}).
    \item[Échantillonnage par attributs] Comptage des anomalies dans un échantillon aléatoire, avec un taux tolérable et une confiance (Ch.~\ref{chap:conduire}).
    \item[Échelle de cotation] Suite de niveaux définis par une description et une plage de valeurs (Ch.~\ref{chap:matrices}).
    \item[Énoncé de risque] Phrase qui relie une cause, un événement et ses conséquences sur un objectif (Ch.~\ref{chap:identifier}).
    \item[Entropie (H)] Dispersion des réponses OUI, NON, NEK, en bits, au plus $\log_2 3$ (Ch.~\ref{chap:retex}).
    \item[Exercice, test] Activité qui entraîne et évalue à partir d'un scénario~; un test attend un résultat mesurable (Ch.~\ref{chap:exercer}).
    \item[FAIR] Méthode d'analyse quantitative~: risque = fréquence des pertes × ampleur des pertes (Ch.~\ref{chap:quantitatif}).
    \item[Fréquence des pertes (LEF)] Nombre d'événements de perte par an~: fréquence des menaces × vulnérabilité (Ch.~\ref{chap:quantitatif}).
    \item[Gouvernance des risques] Rôles, règles et instances qui fixent l'appétence, attribuent et vérifient le traitement des risques (Ch.~\ref{chap:gouvernance}).
    \item[Heuristique] Règle de jugement simple, rapide, source de biais (Ch.~\ref{chap:perception}).
    \item[Incertitude (Knight)] Situation dont on ne connaît pas la distribution des issues (Ch.~\ref{chap:risque}).
    \item[Incertitude (U)] Part des réponses NEK dans une grille (Ch.~\ref{chap:retex}).
    \item[Incident de sécurité] Événements indésirables qui ont une probabilité significative de compromettre l'activité (ISO/IEC~27035) (Ch.~\ref{chap:incidents}).
    \item[Indicateur clé de risque (KRI)] Mesure avancée de l'évolution d'un risque, avec des seuils d'alerte (Ch.~\ref{chap:traiter}).
    \item[Instrument OUI / NON / NEK] Grille dont chaque question admet une troisième réponse~: connaissance insuffisante pour conclure (Ch.~\ref{chap:retex}).
    \item[Limite supérieure] Taux d'anomalies que la population ne dépasse pas, à un niveau de confiance donné (Ch.~\ref{chap:conduire}).
    \item[Main courante] Registre chronologique des faits, informations, décisions et actions de la crise (Ch.~\ref{chap:cellule}).
    \item[Management de crise] Capacité à anticiper une crise, à s'y préparer, à y répondre et à s'en rétablir (Ch.~\ref{chap:crise}).
    \item[Matrice de risques] Tableau vraisemblance × impact qui classe les risques~; elle ne les mesure pas (Cox) (Ch.~\ref{chap:matrices}).
    \item[Menace d'une partie prenante] Dans EBIOS RM, (dépendance × pénétration) / (maturité cyber × confiance) (Ch.~\ref{chap:ebios}).
    \item[Message d'attente] Premier message bref~: ce qui s'est passé, ce que l'on fait, ce que l'on ignore, quand on en dira plus (Ch.~\ref{chap:communiquer}).
    \item[Mode dégradé] Manière de poursuivre une activité à un niveau réduit sans certaines ressources (Ch.~\ref{chap:pca}).
    \item[Modèle des trois lignes] Opérations qui gèrent les risques, fonctions qui aident et surveillent, audit qui donne une assurance (IIA) (Ch.~\ref{chap:gouvernance}).
    \item[Monte-Carlo (simulation)] Tirages aléatoires répétés des facteurs incertains pour étudier la distribution des résultats (Ch.~\ref{chap:simulation}).
    \item[Niveau minimal de continuité (MBCO)] Niveau d'activité acceptable pendant la perturbation (Ch.~\ref{chap:bia}).
    \item[Obligation de moyens] Obligation de mettre en œuvre les moyens raisonnables, sans garantir le résultat (Ch.~\ref{chap:droit}).
    \item[Période d'interruption maximale tolérable (MTPD)] Durée au-delà de laquelle l'impact d'une interruption devient inacceptable (Ch.~\ref{chap:bia}).
    \item[PERT (moyenne)] (minimum + 4 × plus probable + maximum) / 6 (Ch.~\ref{chap:quantitatif}).
    \item[Perte annuelle attendue] Fréquence moyenne × ampleur moyenne~; compare des options, ne décrit pas la pire année (Ch.~\ref{chap:risque}).
    \item[Perte de données maximale admissible (RPO)] Quantité de données, exprimée en durée, que l'on accepte de perdre (Ch.~\ref{chap:bia}).
    \item[Plan de traitement] Actions, justification, responsable, ressources, échéance, indicateur et résiduel visé de chaque risque traité (Ch.~\ref{chap:traiter}).
    \item[Point de situation] Ce que l'on sait, ce que l'on ne sait pas, ce que l'on fait (Ch.~\ref{chap:crise}).
    \item[Preuve d'audit] Enregistrements, énoncés de faits ou informations, pertinents et vérifiables (Ch.~\ref{chap:conduire}).
    \item[Programme d'audit] Ensemble d'audits planifiés pour une durée donnée et dirigés vers un but (Ch.~\ref{chap:audit}).
    \item[Résilience opérationnelle] Capacité à fournir les opérations critiques à travers les perturbations (Ch.~\ref{chap:pra}).
    \item[Retour d'expérience] Démarche qui établit les faits, les écarts, les causes et décide des actions qui changent les pratiques (Ch.~\ref{chap:retex}).
    \item[Risque (ISO~31000)] Effet de l'incertitude sur les objectifs (Ch.~\ref{chap:risque}).
    \item[Risque résiduel] Risque qui subsiste après les seules mesures en place et vérifiées (Ch.~\ref{chap:traiter}).
    \item[Seuil de rentabilité] Valeur d'une hypothèse à partir de laquelle une décision bascule (Ch.~\ref{chap:risque}).
    \item[Signal faible] Information partielle ou ambiguë qui peut annoncer un risque nouveau ou aggravé (Ch.~\ref{chap:identifier}).
    \item[Taxonomie des risques] Répartition des risques en familles stables, pour la couverture et la consolidation (Ch.~\ref{chap:identifier}).
    \item[Théorie de Coombs] Adapter la réponse de communication à la responsabilité attribuée~: victime, accident, évitable (Ch.~\ref{chap:communiquer}).
\end{description}
